
% -----------------------------------------------------------------------------
% Exact scattering foundations before perturbative Feynman rules
% -----------------------------------------------------------------------------

有限时间演化算符描述任意给定时间窗内的动力学；散射算符则进一步要求入射波包在遥远过去、出射波包在遥远未来都已经离开相互作用区。这个渐近条件把“器件内发生了什么”压缩成自由通道之间的映射，也使 Born 近似、Dyson 截断和 Feynman 图都能追溯到同一个精确散射对象。

设
\begin{equation}
 \hat H=\hat H_0+\hat V,
 \label{eq:H-scattering-r5}
\end{equation}
其中 $\hat H_0$ 的绝对连续谱描述远离相互作用区的自由通道。M{\o}ller 波算符定义为
\begin{align}
 \hat\Omega_+
 &\equiv
 \operatorname*{s-lim}_{t\to-\infty}
 \ee^{+\ii\hat Ht/\hbar}
 \ee^{-\ii\hat H_0t/\hbar}
 \hat P_{\rm ac}(\hat H_0),\\
 \hat\Omega_-
 &\equiv
 \operatorname*{s-lim}_{t\to+\infty}
 \ee^{+\ii\hat Ht/\hbar}
 \ee^{-\ii\hat H_0t/\hbar}
 \hat P_{\rm ac}(\hat H_0).
 \label{eq:moller-wave-operators-r5}
\end{align}
若这些强极限存在，则
\begin{equation*}
 |\psi^{(+)}\rangle=\hat\Omega_+|\psi_{\rm in}\rangle,
 \qquad
 |\psi^{(-)}\rangle=\hat\Omega_-|\psi_{\rm out}\rangle,
\end{equation*}
而散射算符定义为
\begin{equation}
 \boxed{\hat S=\hat\Omega_-^\dagger\hat\Omega_+.}
 \label{eq:exact-S-waveop-r5}
\end{equation}
它把遥远过去的自由入射通道映射到遥远未来的自由出射通道。Born 级数、Dyson 展开和 Feynman 图都是求这个精确散射映射的具体展开方法。


定态散射把同一问题改写到固定能量 $E$ 上。设 $|\phi_E\rangle$ 是自由 Hamilton 量 $\hat H_0$ 的广义本征态。由于 $E$ 位于连续谱上，$E-\hat H_0$ 不能作普通有界逆；入射边界条件由预解式边界值
\begin{equation}
 \hat G_0^{(+)}(E)
 \equiv\lim_{\epsilon\downarrow0}
 \frac{1}{E-\hat H_0+\ii\epsilon}
 \label{eq:G0-plus-r5}
\end{equation}
选定。完整入射散射态满足
\begin{equation}
 \boxed{
 |\psi_E^{(+)}\rangle
 =
 |\phi_E\rangle
 +\hat G_0^{(+)}(E)\hat V
 |\psi_E^{(+)}\rangle.}
 \label{eq:LS-state-r5}
\end{equation}
第一项保留指定的自由入射波，第二项是相互作用产生并按出射边界传播的散射波。

把相互作用真正产生的源项定义为
\begin{equation}
 \hat T(E)|\phi_E\rangle
 \equiv
 \hat V|\psi_E^{(+)}\rangle.
 \label{eq:T-def-r5}
\end{equation}
则精确跃迁算符满足
\begin{equation}
 \boxed{
 \hat T(E)
 =
 \hat V+\hat V\hat G_0^{(+)}(E)\hat T(E).}
 \label{eq:T-exact-r5}
\end{equation}
$\hat T(E)$ 因而是把完整散射态压缩到自由通道上的能量分辨算符；Born 级数由这一精确方程对 $\hat V$ 逐次迭代得到。



波算符把自由通道与完整散射态联系起来，同时保持能量生成元的作用：
\begin{equation}
 \boxed{\hat H\hat\Omega_\pm=\hat\Omega_\pm\hat H_0.}
 \label{eq:moller-intertwining-r68}
\end{equation}
因此 $\hat\Omega_\pm|\phi_E\rangle$ 与 $|\phi_E\rangle$ 具有相同连续谱能量。对精确 $T$ 方程逐次迭代，在余项于所研究子空间趋零时得到 Born 展开
\begin{equation}
 \boxed{
 \hat T(E)=\hat V+\hat V\hat G_0^{(+)}\hat V
 +\hat V\hat G_0^{(+)}\hat V\hat G_0^{(+)}\hat V+\cdots.}
 \label{eq:born-series-r68}
\end{equation}
$\hat T$ 的定义由精确 Lippmann--Schwinger 方程给出；Born 近似只发生在级数截断时，其控制量是自由预解式与势的组合，而不是孤立的 $\hat V$。

对时间无关散射，时间平移不变性把散射算符限制在同一能量壳上。采用前文固定的连续谱归一化，$S$ 与能量分辨 $T$ 的矩阵元满足
\begin{equation}
 \boxed{
 \langle f|\hat S|i\rangle
 =\langle f|i\rangle
 -2\pi\ii\,\delta(E_f-E_i)\,
 \langle f|\hat T(E_i)|i\rangle.}
 \label{eq:S-T-energy-r68}
\end{equation}
若散射波算符在所讨论子空间渐近完备，则 $\hat S$ 幺正。写 $\hat S=\id+\ii\hat T_S$ 后可等价改写为
\begin{equation}
 \boxed{
 -\ii(\hat T_S-\hat T_S^\dagger)
 =\hat T_S^\dagger\hat T_S,
 \qquad
 2\Im\langle i|\hat T_S|i\rangle
 =\sum_X|\langle X|\hat T_S|i\rangle|^2.}
 \label{eq:optical-unitarity-r68}
\end{equation}
第二式是在完整在壳出射通道上插入单位算符后的前向矩阵元形式；加入连续相空间测度后即成为通常的光学定理。



外部粒子由 $t\to\pm\infty$ 的渐近条件选成自由哈密顿量的质量壳本征态，因此外部动量满足各自在壳条件。内部传播来自预解式或时间序二点函数，只受各顶角四动量守恒约束，并没有独立的渐近自由边界条件。因此内部四动量一般满足
\begin{equation*}
 q^2\neq m^2c^2.
\end{equation*}
“虚粒子”只是对这类未被外部质量壳条件固定的内部传播变量的简称；它不是一颗在短时间内“违反能量守恒”的真实粒子。


\begin{note}[$S$ 矩阵与有限时间传播]
无限时间散射可用
\begin{equation*}
 \hat V(t)\longrightarrow
 \ee^{-\epsilon|t|}\hat V,
 \qquad
 \epsilon\downarrow0,
\end{equation*}
选择相互作用的真空和边界条件。对有限长度、有限脉冲的自由电子量子光学器件，更基本的对象仍是有限 $t_i,t_f$ 的传播子 $\hat U(t_f,t_i)$；只有当入、出波包在器件外已经形成清楚的自由通道时，再把有限时间传播压缩成 $S$-矩阵才有意义。
\end{note}

\begin{derivation}[从\eqnrefs{eq:moller-wave-operators-r5}到\eqnrefs{eq:moller-intertwining-r68}]

正文\eqnrefs{eq:moller-wave-operators-r5} 用遥远过去和未来的强极限定义 Møller 波算符。把完整时间演化移过这一极限并重命名极限变量，可以证明波算符把 $H_0$ 的时间平移交织成 $H$ 的时间平移；在零时间处求导便得到正文\eqnrefs{eq:moller-intertwining-r68}。



由定义考虑
\begin{align*}
 \ee^{-\ii\hat H\tau/\hbar}\hat\Omega_+
 &=
 \lim_{t\to-\infty}
 \ee^{-\ii\hat H\tau/\hbar}
 \ee^{+\ii\hat Ht/\hbar}
 \ee^{-\ii\hat H_0t/\hbar},\\
 &=
 \lim_{t\to-\infty}
 \ee^{+\ii\hat H(t-\tau)/\hbar}
 \ee^{-\ii\hat H_0(t-\tau)/\hbar}
 \ee^{-\ii\hat H_0\tau/\hbar},\\
 &=\hat\Omega_+\ee^{-\ii\hat H_0\tau/\hbar}.
\end{align*}
在 $\tau=0$ 对两边求导，
\begin{equation*}
 {
 \hat H\hat\Omega_+
 =\hat\Omega_+\hat H_0.}
 \end{equation*}
对 $\hat\Omega_-$ 需要把极限改为 $t\to+\infty$。对任意有限 $\tau$，
\begin{align*}
 \ee^{-\ii\hat H\tau/\hbar}\hat\Omega_-
 &=\lim_{t\to+\infty}
 \ee^{-\ii\hat H\tau/\hbar}
 \ee^{+\ii\hat Ht/\hbar}
 \ee^{-\ii\hat H_0t/\hbar},\\
 &=\lim_{t\to+\infty}
 \ee^{+\ii\hat H(t-\tau)/\hbar}
 \ee^{-\ii\hat H_0(t-\tau)/\hbar}
 \ee^{-\ii\hat H_0\tau/\hbar},\\
 &=\hat\Omega_-\ee^{-\ii\hat H_0\tau/\hbar}.
\end{align*}
在 $\tau=0$ 处分别对两侧求导得到
\[
 -\frac{\ii}{\hbar}\hat H\hat\Omega_-
 =-\frac{\ii}{\hbar}\hat\Omega_-\hat H_0,
\]
因此
\[
 \hat H\hat\Omega_-=\hat\Omega_-\hat H_0.
\]
于是若 $|\phi_E\rangle$ 是 $\hat H_0$ 的广义能量本征态，
\[
 \hat H\,\hat\Omega_\pm|\phi_E\rangle
 =\hat\Omega_\pm\hat H_0|\phi_E\rangle
 =E\,\hat\Omega_\pm|\phi_E\rangle,
\]
所以 $\hat\Omega_\pm|\phi_E\rangle$ 是完整哈密顿量在同一连续能量上的入、出散射态。
\end{derivation}
\begin{derivation}[从\eqnrefs{eq:G0-plus-r5}到\eqnrefs{eq:LS-state-r5}]

连续谱能量 $E$ 上不能把 $E-H_0$ 当作普通有界算符直接求逆，因此正文\eqnrefs{eq:G0-plus-r5} 先用 $+\ii0$ 选定出射预解式边界值。把定态 Schrödinger 方程移项后作用这一边界值，并保留自由齐次解，就得到正文\eqnrefs{eq:LS-state-r5}。



令
\begin{equation*}
 (\hat H_0+\hat V)|\psi_E^{(+)}\rangle
 =E|\psi_E^{(+)}\rangle.
\end{equation*}
移项得到
\begin{equation*}
 (E-\hat H_0)|\psi_E^{(+)}\rangle
 =\hat V|\psi_E^{(+)}\rangle.
\end{equation*}
因为 $E$ 位于 $\hat H_0$ 的连续谱上，$E-\hat H_0$ 不能作为普通有界算符直接求逆；入射边界条件由预解式的边界值
\begin{equation*}
 \hat G_0^{(+)}(E)
 \equiv
 \lim_{\epsilon\downarrow0}
 \frac{1}{E-\hat H_0+\ii\epsilon}
 \end{equation*}
选定。于是
\begin{equation*}
 {
 |\psi_E^{(+)}\rangle
 =
 |\phi_E\rangle
 +\hat G_0^{(+)}(E)\hat V
 |\psi_E^{(+)}\rangle,}
 \end{equation*}
其中 $(E-\hat H_0)|\phi_E\rangle=0$。若要求出射边界条件，则改用
\begin{equation*}
 \hat G_0^{(-)}(E)
 =
 \lim_{\epsilon\downarrow0}
 \frac{1}{E-\hat H_0-\ii\epsilon}.
\end{equation*}
所以散射边界条件与 Green 函数在连续谱两侧的解析边界值是同一结构。
\end{derivation}
\begin{derivation}[从\eqnrefs{eq:T-exact-r5}到\eqnrefs{eq:born-series-r68}]

正文\eqnrefs{eq:T-exact-r5} 已经是完整散射的精确算符方程。将右端的 $T$ 连续代回自身，可把一次、两次、三次散射按 $V G_0^{(+)}$ 的次数排列；只有当剩余项在所研究子空间上可忽略时，有限恒等式才变成正文\eqnrefs{eq:born-series-r68} 的 Born 展开。


从精确方程
\begin{equation*}
T=V+VG_0^{(+)}T
\end{equation*}
出发。第一次把右端的 $T$ 再代回：
\begin{align*}
T
&=V+VG_0^{(+)}\big(V+VG_0^{(+)}T\big)\\
&=V+VG_0^{(+)}V+VG_0^{(+)}VG_0^{(+)}T.
\end{align*}
再迭代一次最后的 $T$：
\begin{align*}
T
={}&V+VG_0^{(+)}V+VG_0^{(+)}VG_0^{(+)}V\\
&+VG_0^{(+)}VG_0^{(+)}VG_0^{(+)}T.
\end{align*}
重复这一过程 $N$ 次得到有限阶恒等式
\begin{equation*}
T=\sum_{n=0}^{N}(VG_0^{(+)})^nV+(VG_0^{(+)})^{N+1}T.
\end{equation*}
若余项在所研究初末态或工作子空间上趋零，才可写成 Born 级数
\begin{equation*}
T=V+VG_0^{(+)}V+VG_0^{(+)}VG_0^{(+)}V+\cdots.
\end{equation*}
因此把一阶、二阶 Born 截断分别定义为
\begin{equation*}
T^{[1]}\equiv V,
\qquad
T^{[2]}\equiv V+VG_0^{(+)}V.
\end{equation*}
上面的有限阶恒等式给出相应精确余项
\begin{equation*}
T-T^{[1]}=VG_0^{(+)}T,
\qquad
T-T^{[2]}=(VG_0^{(+)})^2T.
\end{equation*}
只有这些余项在所研究矩阵元或工作子空间范数中足够小时，$T^{[1]}$ 或 $T^{[2]}$ 才是受控近似。无量纲控制对象因此是“自由传播一次以后再散射”的组合，例如投影到工作子空间后的 $\|G_0^{(+)}V\|$ 或相应矩阵元。
\end{derivation}
\begin{derivation}[从\eqnrefs{eq:exact-S-waveop-r5}到\eqnrefs{eq:S-T-energy-r68}]
这里要分清两件事：\(\delta(E_f-E_i)\) 来自无限时间平移对称，而 \(T(E_i)\) 包含的是作用期间所有重复散射。二者可以从精确相互作用绘景积分方程中直接分离。

设
\begin{equation*}
\hat H_0|i\rangle=E_i|i\rangle,
\qquad
\hat H_0|f\rangle=E_f|f\rangle,
\end{equation*}
并令入射定态散射态
\begin{equation*}
|\psi_i^{(+)}\rangle
\equiv\hat\Omega_+|i\rangle.
\end{equation*}
由前一推导的交织关系
\(\hat H\hat\Omega_+=\hat\Omega_+\hat H_0\)，有
\begin{equation*}
\hat H|\psi_i^{(+)}\rangle
=E_i|\psi_i^{(+)}\rangle.
\end{equation*}
因此该态在 Schr\"odinger 绘景只积累定态相位，而其相互作用绘景表示为
\begin{align*}
|\psi_{i,I}^{(+)}(t)\rangle
&=\ee^{+\ii\hat H_0t/\hbar}
\ee^{-\ii\hat Ht/\hbar}|\psi_i^{(+)}\rangle\\
&=\ee^{+\ii\hat H_0t/\hbar}
\ee^{-\ii E_it/\hbar}|\psi_i^{(+)}\rangle.
\end{align*}

精确相互作用绘景传播子满足 Volterra 方程
\begin{equation*}
\hat U_I(+\infty,-\infty)
=\id-\frac{\ii}{\hbar}
\int_{-\infty}^{+\infty}\dd t\,
\hat V_I(t)\hat U_I(t,-\infty).
\end{equation*}
对入射态 \(|i\rangle\)，绝热极限下
\(\hat U_I(t,-\infty)|i\rangle=|\psi_{i,I}^{(+)}(t)\rangle\)。因此
\begin{align*}
\langle f|\hat S|i\rangle
={}&\langle f|i\rangle
-\frac{\ii}{\hbar}
\int_{-\infty}^{+\infty}\dd t\,
\langle f|\hat V_I(t)|\psi_{i,I}^{(+)}(t)\rangle.
\end{align*}
把
\(\hat V_I(t)=\ee^{+\ii\hat H_0t/\hbar}\hat V\ee^{-\ii\hat H_0t/\hbar}\)
和上面的定态形式代入，\(\hat H_0\) 在 bra 上给出 \(E_f\)，于是时间依赖完全提出：
\begin{align*}
\langle f|\hat V_I(t)|\psi_{i,I}^{(+)}(t)\rangle
&=\ee^{+\ii(E_f-E_i)t/\hbar}
\langle f|\hat V|\psi_i^{(+)}\rangle.
\end{align*}

再由 Lippmann--Schwinger 方程
\begin{equation*}
|\psi_i^{(+)}\rangle
=|i\rangle+\hat G_0^{(+)}(E_i)\hat V|\psi_i^{(+)}\rangle
\end{equation*}
以及正文对 \(\hat T(E)\) 的定义，可得
\begin{equation*}
\hat V|\psi_i^{(+)}\rangle
=\hat T(E_i)|i\rangle.
\end{equation*}
这一等式把任意次数的
\(V G_0^{(+)}V G_0^{(+)}\cdots V\)
全部收进了 \(T(E_i)\)，剩余时间积分于是给出整体能量守恒：
\begin{equation*}
\int_{-\infty}^{+\infty}\dd t\,
\ee^{+\ii(E_f-E_i)t/\hbar}
=2\pi\hbar\,\delta(E_f-E_i).
\end{equation*}
于是
\begin{equation*}
\langle f|\hat S|i\rangle
=\langle f|i\rangle
-2\pi\ii\,\delta(E_f-E_i)
\langle f|\hat T(E_i)|i\rangle,
\end{equation*}
即正文\eqnrefs{eq:S-T-energy-r68}。

这个推导同时说明了两个常被混淆的对象：\(\delta(E_f-E_i)\) 只反映时间平移不变性；\(T(E_i)\) 才承载相互作用区内全部多次散射动力学。相对论局域场论再把空间平移对称性并入同一逻辑，得到四维动量守恒 \(\delta^{(4)}\)。
\end{derivation}
\begin{derivation}[从\eqnrefs{eq:exact-S-waveop-r5}到\eqnrefs{eq:optical-unitarity-r68}]

正文\eqnrefs{eq:exact-S-waveop-r5} 在渐近完备的散射子空间上由两个波算符组成，因此保持总概率。把 $S$ 写成单位算符与跃迁部分之和，展开 $S^\dagger S=\id$，再在前向入射态之间插入完整出射通道，就得到正文\eqnrefs{eq:optical-unitarity-r68}。



若波算符渐近完备，则散射子空间上
\begin{equation*}
 \hat S^\dagger\hat S=\id.
\end{equation*}
写
\begin{equation*}
 \hat S=\id+\ii\hat T_S,
\end{equation*}
其中 $\hat T_S$ 是无量纲跃迁算符。代入得
\begin{align*}
 (\id-\ii\hat T_S^\dagger)
 (\id+\ii\hat T_S)
 &=\id,\\
 \ii(\hat T_S-\hat T_S^\dagger)
 +\hat T_S^\dagger\hat T_S
 &=0.
\end{align*}
所以
\begin{equation*}
 {
 -\ii(
 \hat T_S-\hat T_S^\dagger)
 =
 \hat T_S^\dagger\hat T_S.}
 \end{equation*}
在两侧夹入同一入射态并插入完整的在壳出射通道：
\begin{equation*}
 2\,\operatorname{Im}
 \langle i|\hat T_S|i\rangle
 =
 \sum_X
 |\langle X|\hat T_S|i\rangle|^2.
 \end{equation*}
右端是所有开放通道的正概率之和。对二体入射，将连续态求和写成 Lorentz 不变相空间后，便得到光学定理：前向振幅的虚部与总截面成正比。
\end{derivation}
